%% ---------------------------------------------------------------------
%%  I Radicali: la Scheda Unica — campo di esistenza, operazioni e
%%  semplificazione in UNA pagina A4.
%%
%%  Ordine dell'operatore, 26/11/2026: «la scheda delle formule deve essere
%%  unita alla scheda delle semplificazioni: deve essere un'unica scheda con
%%  tutte le formule». Prima erano due file (radicali-1 e radicali-2): qui
%%  c'è tutto, e le sezioni sono le tre di prima nell'ordine in cui si
%%  studiano — prima DOVE la radice esiste, poi come si opera, poi come si
%%  semplifica.
%%
%%  DUE CORREZIONI ALLA SCHEDA (ordine dell'operatore, 27/11/2026):
%%  (1) «portar dentro» mette il fattore a DESTRA di quello che era già
%%      sotto radice — ²√(▲ · ●²), non ²√(●² · ▲);
%%  (2) la semplificazione si legge al CONTRARIO — «cambiare l'indice della
%%      radice»: ³√(▲) = ⁶√(▲²), cioè l'indice si moltiplica, non si divide.
%%  Le altre sezioni, la libreria e la palette non si toccano.
%%
%%  Qui non c'è grafica: tutto sta in ../schede.sty. Le formule sono quelle
%%  della scheda dell'operatore (il PDF Chromium «Schema sui Radicali»,
%%  124364 B, in tests/golden/schede/) e quelle della seconda scheda: nessuna
%%  è stata tolta per far stare la pagina — il respiro sì, ed è la libreria a
%%  deciderlo (\schedaCompatta, \schedaMarginiSicuri).
%% ---------------------------------------------------------------------
\documentclass[10pt]{article}
\usepackage{schede}

%% Testata e piede dentro la zona che ogni stampante sa stampare (senza,
%% il piede finiva su un secondo foglio: misurato il 26/11).
\schedaMarginiSicuri
%% Tre sezioni e undici righe in una pagina: si toglie respiro agli SPAZI
%% (righe più basse, margini delle celle più stretti), non al testo.
\schedaCompatta

\schedaTitolo{I Radicali: la Scheda Unica}
             {CE \textperiodcentered\ Operazioni \textperiodcentered\ Semplificazione}
\schedaPiede{Schema di Matematica — I Radicali}

\begin{document}

%% =====================================================================
\schedaSezione{1}{Campo di Esistenza (CE)}
%% =====================================================================
\begin{schedaTabella}{Tipo di radice}{Schema visivo}{Esempio}

  \riga{\voce{Indice pari}{(es. $\sqrt{\ }$, $\sqrt[4]{\ }$, $\sqrt[6]{\ }$)}}
       {\cePari{\formaQuadrato}}
       {\cePari[4]{\formaQuadrato}}

  \riga{\voce{Indice dispari}{(es. $\sqrt[3]{\ }$, $\sqrt[5]{\ }$, $\sqrt[7]{\ }$)}}
       {\ceDispari[3]{\formaQuadrato}}
       {\ceDispari[5]{\formaQuadrato}}

  \riga{\voce{Più radici pari}{(sistema di condizioni)}}
       {$\radice{\formaQuadrato}+\radice{\formaTriangolo}\;\freccia\;
         \sistema{\formaQuadrato\ge 0\\ \formaTriangolo\ge 0}$}
       {$\radice{\formaQuadrato}+\radice[4]{\formaTriangolo}\;\freccia\;
         \sistema{\formaQuadrato\ge 0\\ \formaTriangolo\ge 0}$}

\end{schedaTabella}

%% =====================================================================
\schedaSezione{2}{Operazioni con i Radicali (Casi Semplici)}
%% =====================================================================
\begin{schedaTabella}{Operazione}{Schema visivo (stesso indice $n$)}{Esempio}

  \riga{\voce{Somma e sottrazione}{(solo radicali simili)}}
       {$\mathbf{a}\,\radice[n]{\formaQuadrato}\;\pm\;\mathbf{b}\,\radice[n]{\formaQuadrato}
         \;\freccia\;(a\pm b)\,\radice[n]{\formaQuadrato}$\par\vspace{3pt}
        {\footnotesize\color{schedeGrigio}Con radici diverse:
         $a\radice{\formaQuadrato}+b\radice{\formaTriangolo}$ resta così}}
       {$3\radice{\formaQuadrato}+5\radice{\formaQuadrato}
         =8\radice{\formaQuadrato}$\par\vspace{3pt}
        $4\radice{\formaQuadrato}+2\radice{\formaTriangolo}
         -\radice{\formaQuadrato}$\par
        $=3\radice{\formaQuadrato}+2\radice{\formaTriangolo}$}

  \riga{\voce{Prodotto}{(moltiplicazione)}}
       {$\radice[n]{\formaQuadrato}\cdot\radice[n]{\formaTriangolo}
         \;\freccia\;\radice[n]{\formaQuadrato\cdot\formaTriangolo}$}
       {$\radice[3]{\formaQuadrato}\cdot\radice[3]{\formaTriangolo}
         =\radice[3]{\formaQuadrato\cdot\formaTriangolo}$}

  \riga{\voce{Divisione}{(quoziente)}}
       {$\radice[n]{\formaQuadrato}:\radice[n]{\formaTriangolo}
         \;\freccia\;\radice[n]{\formaQuadrato:\formaTriangolo}$\par\vspace{3pt}
        {\footnotesize\color{schedeGrigio}(con $\formaTriangolo\neq 0$)}}
       {$\radice[4]{\formaQuadrato}:\radice[4]{\formaTriangolo}
         =\radice[4]{\formaQuadrato:\formaTriangolo}$}

  \riga{\voce{Potenza di radice}{}}
       {$(\radice[n]{\formaQuadrato})^{k}
         \;\freccia\;\radice[n]{\formaQuadrato^{k}}$}
       {$(\radice[3]{\formaQuadrato})^{2}
         =\radice[3]{\formaQuadrato^{2}}$\par\vspace{3pt}
        $(2\radice[3]{\formaQuadrato})^{2}
         =4\radice[3]{\formaQuadrato^{2}}$}

  \riga{\voce{Radice di radice}{}}
       {$\radice[m]{\radice[n]{\formaQuadrato}}
         \;\freccia\;\radice[m\cdot n]{\formaQuadrato}$}
       {$\radice[3]{\radice[2]{\formaQuadrato}}
         =\radice[3\cdot 2]{\formaQuadrato}
         =\radice[6]{\formaQuadrato}$}

\end{schedaTabella}

%% =====================================================================
\schedaSezione{3}{Semplificazione: la Proprietà Invariantiva}
%% =====================================================================

\riquadro[La regola]{Moltiplicare o dividere per lo stesso numero (diverso
da zero) l'indice della radice e l'esponente del radicando non cambia la
radice: è la stessa radice scritta in un altro modo.}

\vspace{6pt}

\begin{schedaTabella}{Cosa si fa}{Schema visivo (stesso numero $n$)}{Esempio}

  \riga{\voce{Cambiare l'indice della radice}{(moltiplicare per $k$)}}
       {$\radice[n]{\formaQuadrato}
         \;\freccia\;\radice[n\cdot k]{\formaQuadrato^{k}}$}
       {$\radice[3]{\formaQuadrato}
         =\radice[6]{\formaQuadrato^{2}}$}

  \riga{\voce{Portar fuori}{(se il radicando è un prodotto)}}
       {$\radice[n]{\formaQuadrato^{n}\cdot\formaTriangolo}
         \;\freccia\;\formaQuadrato\cdot\radice[n]{\formaTriangolo}$}
       {$\radice[3]{\formaQuadrato^{3}\cdot\formaTriangolo}
         =\formaQuadrato\cdot\radice[3]{\formaTriangolo}$}

  \riga{\voce{Portar dentro}{(il fattore davanti)}}
       {$\formaQuadrato\cdot\radice[n]{\formaTriangolo}
         \;\freccia\;\radice[n]{\formaTriangolo\cdot\formaQuadrato^{n}}$}
       {$\formaCerchio\cdot\radice[2]{\formaTriangolo}
         =\radice[2]{\formaTriangolo\cdot\formaCerchio^{2}}$}

\end{schedaTabella}

{\footnotesize\color{schedeGrigio}Portare fuori e portare dentro sono la
stessa uguaglianza; cambiare l'indice è la proprietà invariantiva letta
dall'altro verso.\par}

\end{document}
